%% ma: L12-B13-0015
%% lop: 12
%% chuong: 4
%% bai: 13
%% dang: 12.13.8
%% loai: TLN
%% muc: VD
%% dap_an: "1,26"
%% nguon: "(NBV)-ĐỀ SỐ 5-MỖI NGÀY 1 ĐỀ THI 2025.pdf (D05, MỖI NGÀY 1 ĐỀ - Đề số 5), Phần III Câu 3"
%% kien_thuc: "Bài 13 (diện tích hình phẳng giới hạn bởi parabol, tính bằng tích phân); Bài 16 lớp 10 (hàm số bậc hai)"
%% trang_thai: cho-duyet
%% ly_do: "Dạng toán mới đề xuất 12.13.8: chia cổng parabol thành các phần có diện tích bằng nhau, tìm tỉ số hai dây"
%% ghi_chu: "Mức VD (phân vân giữa VD và VDC, chọn thấp hơn). Vị trí hai dây trên hình vẽ minh họa (tỉ số CD/AB không phụ thuộc vị trí dây AB); đã vẽ đúng tỉ lệ s/t = căn bậc ba của 2"
\begin{ex}
Một cổng có dạng hình parabol với chiều cao $8$ m, chiều rộng chân đế $8$ m (hình vẽ). Người ta căng hai sợi dây trang trí $AB$, $CD$ nằm ngang, đồng thời chia cổng thành ba phần sao cho hai phần ở phía trên có diện tích bằng nhau. Tỉ số $\dfrac{CD}{AB}$ bằng bao nhiêu (làm tròn kết quả đến hàng phần trăm)?
\begin{center}
\begin{tikzpicture}[x=.55cm,y=.4cm]
  \draw[thick,domain=-4:4,samples=70,smooth] plot (\x,{8-\x*\x/2});
  \draw (-4,0)--(4,0);
  \draw (-2.2,5.58)--(2.2,5.58) (-2.772,4.158)--(2.772,4.158);
  \draw[<->] (-4,-1)--(4,-1) node[midway,below]{$8$ m};
  \draw[<->] (5.2,0)--(5.2,8) node[midway,right]{$8$ m};
  \draw[phu] (4,0)--(5.5,0) (0,8)--(5.5,8);
  \path (-2.2,5.58) node[above left]{$A$} (2.2,5.58) node[above right]{$B$}
        (-2.772,4.158) node[left]{$C$} (2.772,4.158) node[right]{$D$};
\end{tikzpicture}
\end{center}
\shortans{1,26}
\loigiai{Chọn hệ trục với gốc ở chân đế, trục $Oy$ qua đỉnh: parabol có đỉnh $(0;8)$ và đi qua $(4;0)$ nên $y=8-\dfrac{x^2}{2}$.
Phần parabol nằm phía trên dây ngang tại hoành độ $\pm t$ ($t>0$) có diện tích $\displaystyle\int_{-t}^{t}\left[\left(8-\dfrac{x^2}{2}\right)-\left(8-\dfrac{t^2}{2}\right)\right]\mathrm dx=\dfrac23t^3$.
Đặt $AB=2t$, $CD=2s$. Phần trên $AB$ có diện tích $\dfrac23t^3$; phần giữa $AB$ và $CD$ có diện tích $\dfrac23s^3-\dfrac23t^3$. Hai phần bằng nhau nên $s^3=2t^3$, suy ra $\dfrac{CD}{AB}=\dfrac st=\sqrt[3]2\approx1{,}26$.}
\end{ex}
